\section{Plantillas}\label{sec:plantillas}

Muchas clases repiten las mismas formas: una definición, un suceso, un proceso,
una comparación, un argumento, un estudio. Estas plantillas son esqueletos para
esas situaciones. Las uso cuando encajan y las dejo cuando no.

\subsection{Concepto}
Definición, ejemplo, con qué no confundirlo y a qué tipo pertenece.
\begin{apunte}[Plantilla · concepto]
\linea{}{término \ks{K09} definición breve}
\linea{}{término \(\subset\) tipo más general}
\linea{}{\ks{K26} caso concreto}
\linea{}{término \(\neq\) concepto parecido: la diferencia}
\end{apunte}
\begin{apunte}[Química · ejemplo]
\linea{}{catalizador \ks{K09} \(\uparrow\) velocidad de reacción s/ consumirse}
\linea{}{enzima \(\subset\) catalizador}
\linea{}{\ks{K26} amilasa: almidón \ks{K24} azúcares}
\linea{\ks{K31}}{catalizador \(\neq\) reactivo: no aparece en los productos}
\end{apunte}

\subsection{Suceso}
Para historia, política, noticias o biografías.
\begin{apunte}[Plantilla · suceso]
\linea{}{\ks{K07} cuándo \quad \ks{K08} dónde \quad \ks{K30} quiénes}
\linea{}{\ks{K06} causas \ks{K23} suceso}
\linea{}{suceso \ks{K23} \ks{K06} consecuencias}
\linea{}{\ks{K16} fuente}
\end{apunte}
\begin{apunte}[Historia · ejemplo]
\linea{}{Crac del 29 \quad \ks{K07} oct.~1929 \quad \ks{K08} Nueva York}
\linea{}{[especulacn \(\wedge\) crédito fácil] \ks{K23} burbuja bursátil \ks{K21} crac}
\linea{}{crac \ks{K23} quiebra de bancos \ks{K23} \(\uparrow\) desempleo}
\linea{\ks{K35}}{crac \ks{K21} Gran Depresión (\ks{K07} 1929--1939)}
\end{apunte}

\subsection{Proceso o algoritmo}
Una etapa por línea, unidas con Luego. Si una etapa produce la siguiente, uso
Causa o Se vuelve en su lugar.
\begin{apunte}[Plantilla · proceso]
\linea{}{\ks{K04} nombre del proceso: entrada \ks{K24} salida}
\linea{}{\ks{K21} etapa 1}
\linea{}{\ks{K21} etapa 2}
\linea{}{\ks{K21} etapa 3 \ks{K10}\(_{\text{fig.}}\)}
\end{apunte}
\begin{apunte}[Computación · ejemplo]
\linea{}{\ks{K04} compilacn: código fuente \ks{K24} ejecutable}
\linea{}{\ks{K21} análisis léxico: texto \ks{K24} tokens}
\linea{}{\ks{K21} análisis sintáctico: tokens \ks{K24} árbol}
\linea{}{\ks{K21} generacn de código}
\linea{\ks{K19}}{¿dónde entra la optimizacn?}
\end{apunte}

\subsection{Comparación}
Una tabla de dos columnas; al final, lo que comparten y lo que los separa.
\begin{apunte}[Filosofía · ejemplo]
\linea{}{\ks{K30} Platón \ks{K25} \ks{K30} Aristóteles}
\linea{}{ideas separadas del mundo \ks{K25} formas en las cosas}
\linea{}{conocer \(\approx\) recordar \ks{K25} conocer \(\approx\) observar}
\linea{\ks{K35}}{\(\approx\): buscan lo universal; \(\neq\): dónde está}
\end{apunte}

\subsection{Argumento}
Premisas numeradas y conclusión con \(\therefore\). Mis objeciones van marcadas
como Idea mía para no confundirlas con el argumento del autor.
\begin{apunte}[Filosofía · ejemplo]
\linea{}{\ks{K16} Descartes:}
\linea{}{1. dudo de todo}
\linea{}{2. dudar \(\to\) pensar}
\linea{}{3. pensar \(\to\) existir}
\linea{\ks{K31}}{\(\therefore\) existo}
\linea{\ks{K32}}{¿2 supone ya un «yo» q piensa?}
\end{apunte}

\subsection{Estudio o experimento}
Hipótesis, qué se midió, qué se encontró y por qué hay que tener cuidado.
\begin{apunte}[Psicología · ejemplo]
\linea{\ks{K33}}{recompensa inmediata \ks{K23} \(\downarrow\) autocontrol}
\linea{}{\ks{K16} Mischel: niños \(\wedge\) malvavisco, \ks{K05} 15 min de espera}
\linea{}{esperar \ks{K28} mejores resultados escolares}
\linea{\ks{K19}}{¿clase social \ks{K23} ambas?}
\end{apunte}

\subsection{Clasificación}
Un árbol con \(\subset\): cada nivel es un tipo del anterior.
\begin{apunte}[Biología · ejemplo]
\linea{}{vertebrados}
\linea{}{\quad peces \(\subset\) vertebrados}
\linea{}{\quad mamíferos \(\subset\) vertebrados}
\linea{}{\qquad primates \(\subset\) mamíferos \quad \ks{K26} chimpancé}
\end{apunte}
